\documentclass[11pt]{article}
\usepackage[margin=1in]{geometry}
\usepackage[T1]{fontenc}
\usepackage{lmodern,color}
\usepackage{parskip}
\usepackage{enumitem}
\usepackage{booktabs}
\usepackage{tabularx}
\usepackage{array}
\usepackage[hidelinks]{hyperref}
\usepackage{xurl}
\newcolumntype{Y}{>{\raggedright\arraybackslash}X}
\newcolumntype{P}[1]{>{\raggedright\arraybackslash}p{#1}}
\newcolumntype{C}[1]{>{\centering\arraybackslash}p{#1}}
\setlist{leftmargin=*,topsep=2pt,itemsep=1pt}
\setlength{\parskip}{5pt}
\setlength{\parindent}{0pt}
\newcommand{\syllabussectionspace}{\vspace{-3mm}}
\makeatletter
\renewcommand\section{\@startsection{section}{1}{0pt}%
  {-1.5ex plus -.4ex minus -.2ex}%
  {.5ex plus .2ex}%
  {\normalfont\Large\bfseries}}
\makeatother
\begin{document}

\begin{center}{\large\textbf{MA213 -- BASIC STATISTICS AND PROBABILITY -- FALL 2026}}\\[2mm]
SYLLABUS\end{center}

\textbf{LECTURES:} Monday, Wednesday, Friday, 11:15AM--12:05PM, PHO 206

\textbf{DISCUSSION AND LAB SECTIONS:} Check your schedule.

\textbf{LECTURE INSTRUCTOR:}
\begin{itemize}
\item Dr. Yongho Lim (\href{mailto:ylim2@bu.edu}{ylim2@bu.edu}) --- Office: CCDS 408, Office hours: Monday 2:00PM-3:30PM; Friday 12:30PM-1:30PM
\end{itemize}
\textbf{LAB INSTRUCTOR:}
\begin{itemize}
\item Dr. Emily Stephen (\href{mailto:estephen@bu.edu}{estephen@bu.edu}) --- Office: CCDS 437, Office hours: Monday 12:15PM-1:30PM; Thursday 2:00PM-3:00PM
\end{itemize}
\textbf{TEACHING FELLOWS:}
\begin{itemize}
\item Lucia Vilallonga (\href{mailto:lvilallonga@bu.edu}{lvilallonga@bu.edu}) --- Office: CCDS 542, Office hours: Tuesday 2:00PM-3:00PM
\item Noah Kupinsky (\href{mailto:n0ah@bu.edu}{n0ah@bu.edu}) --- Office: CCDS 326, Office hours: Tuesday 9:30AM-10:45AM
\item Sydney Elliot (\href{mailto:sydelli@bu.edu}{sydelli@bu.edu}) --- Office: Zoom \url{https://bostonu.zoom.us/j/97322041821?pwd=uzB2ApmMgblECb3GObaX4Xw1aIMRtA.1}, Office hours: Wednesday 2:30PM-3:30PM
\end{itemize}
\textbf{LEARNING ASSISTANTS:}
\begin{itemize}
\item Carly Kent (\href{mailto:kentcar@bu.edu}{kentcar@bu.edu}) --- Office: Tbd, Office hours: Tuesday 11:00AM-12:00PM
\item Jacob Freedman (\href{mailto:jacobf@bu.edu}{jacobf@bu.edu}) --- Office: Tbd, Office hours: Wednesday 1:00PM-2:00PM
\end{itemize}

\section*{}
\textbf{COURSE DESCRIPTION:}

This course serves as an introduction to basic concepts and tools in probability and statistics. We begin with techniques for describing data. Then we study the elements of probability theory, including probability densities, means, variances, correlation, independence, the binomial distribution, and the central limit theorem. Finally, we combine data description and probability theory into an approach to statistical inference. Students should emerge from this course with the ability to incorporate a variety of skills in analyzing and reasoning from data.

\section*{}
\textbf{COURSE MATERIALS:}

\textbf{TEXTBOOK: OpenIntro Statistics, 4th Edition} (required): Open-source textbook by Ce\-tin\-kaya-\allowbreak Run\-del, Diez, and Barr. A PDF is available at no cost; optional print copies are available from the publisher.

\begin{center}\url{https://www.openintro.org/book/os/}\end{center}

\textbf{Edfinity} (required): Used for in-class activities and weekly auto-graded homework. The Fall 2026 reference price was \$40 with OfficeHours.

\section*{}
\textbf{EDFINITY ENROLLMENT:}
\begin{samepage}
\begin{enumerate}
\item[1.] Use a current version of Google Chrome or Firefox on a Windows or Mac computer. Other browsers or devices may cause issues when Edfinity is accessed through Blackboard.
\item[2.] If you already have an Edfinity account, sign out of Edfinity before beginning enrollment.
\item[3.] Log in to Blackboard (learn.bu.edu).
\item[4.] Open the Edfinity link for this course in Blackboard.
\item[5.] Follow the prompts to enroll and pay the course fee.
\end{enumerate}
\end{samepage}

\textbf{Laptop, phone, or similar device} (required): Bring a device to lectures to participate in in-class activities.

\textbf{Pen or pencil} (required): Bring a pen or pencil to lectures and labs to complete the worksheets.

\textbf{R and RStudio} (required): Used during skills labs for statistical programming and data analysis.

\begin{center}\url{https://posit.co/download/rstudio-desktop/}\end{center}

\textbf{Gradescope} (required): Used to submit lab deliverables and revisions.

\section*{}
\textbf{ASSESSMENTS AND GRADING:}

This course uses an alternative grading system based on the following principles:

\begin{samepage}
\begin{itemize}
\item You should learn the course material in a way that promotes sense-making and retention of knowledge.
\item You should spend much of your time working at your ``learning edge'', where you are challenged but not lost or overwhelmed.
\item There should be clear expectations for what quality work looks like and how grades are determined.
\item You should get feedback that helps you improve your work.
\item You should have many opportunities to learn from mistakes that don’t negatively impact your grade.
\item Frequent small assignments/assessments support learning more than infrequent high stakes assignments/assessments do.
\item The class should be a learning community, where students support each other to increase everyone’s learning. You should not be competing with other students for grades.
\item The assessment system should promote authentic engagement with statistics rather than “studenting” behaviors, such as cramming, cheating, and arguing about points.
\item You should have some choice in your trajectory through the course.
\end{itemize}
\end{samepage}

You will complete assignments that fall into different categories: 12 Edfinity Homeworks, in-Lecture worksheet, 5 Quizzes, 7 Skills Labs, and 2 Lab Group Projects. Quiz problems are divided into 18 core and 11 auxiliary learning objectives. Your course grade will be determined by how many items you’ve passed in each category; \textbf{You must meet all of the criteria for a particular grade} (see table below). You will have an opportunity to retake quiz problems and revise projects (see below). You can choose what grade you would like to aim for, and you only need to do the work for that grade (I encourage everyone to aim for an A!). The following table and the accompanying notes summarize the requirements for each grade:

\section*{}
\textbf{GRADE REQUIREMENTS:}

\begin{center}
\scriptsize
\setlength{\tabcolsep}{4pt}
\begin{tabularx}{\textwidth}{lYYYYYYY}
\toprule
\textbf{Grade} & \textbf{Homeworks} & \textbf{Quiz Core Objectives} & \textbf{Quiz Auxiliary Objectives} & \textbf{Skills Labs} & \textbf{Lab Projects} & \textbf{Lectures} & \textbf{Discussions} \\
\midrule
A & 12/12 complete & 16/18 passed & 9/11 passed & 7/7 complete & 2 satisfactory & Fewer than 5 missed & Fewer than 3 missed \\
B & 11/12 complete & 14/18 passed & 6/11 passed & Labs 1-6 complete & 2 satisfactory & Fewer than 9 missed & Fewer than 4 missed \\
C & 10/12 complete & 10/18 passed & 0/11 passed & Labs 1-6 complete & 2 satisfactory & Fewer than 15 missed & Fewer than 5 missed \\
\bottomrule
\end{tabularx}
\end{center}

\section*{}
\textbf{ADDITIONAL FACTORS:}

\begin{samepage}
\begin{itemize}
\item Grades between letter grades are determined by how close a student is to the next letter grade.
\item One unsatisfactory project lowers the course grade by one-third of a letter grade, such as B to B-.
\item Two unsatisfactory project grades lower the course grade by two-thirds of a letter grade, such as B to C+.
\item Lectures/Discussion missed with permission are not counted as missed. To request permission, contact the lecture instructor (for lectures) or discussion TF (for discussions) in advance with the reason for the absence and complete any required form.
\item If you need to miss a lab, contact the lab instructor in advance to get instructions on how to make up the missed work.
\end{itemize}
\end{samepage}

\section*{}
\textbf{HOMEWORK:}

Homework problem sets will be assigned in Edfinity each week and due the following Monday at 3pm. The interface will tell you immediately if your answer is incorrect, and you can try multiple times. In order to get credit for a complete homework, you must get at least 80\% of the problems correct.

\section*{}
\textbf{QUIZZES:}

The course material is divided into a set of Learning Objectives, described in the Learning Objectives document on the course website. There are 18 Core objectives and 11 Auxiliary objectives that will be assessed via 5 quizzes that occur during lecture time. Note that Quiz 5 is entirely Auxiliary objectives. On any given quiz, each objective will be graded Pass/Almost Pass/Not Yet.

\begin{center}
\small
\setlength{\tabcolsep}{4pt}
\begin{tabularx}{\textwidth}{lYY}
\toprule
\textbf{Result} & \textbf{Meaning} & \textbf{Next Step} \\
\midrule
Pass & The response is correct. & The objective does not need to be attempted again, even when another objective from the same quiz is retaken. \\
Almost Pass & The response contains only a minor error, such as a typo. & Within one week of receiving the quiz grade, attend office hours, explain the error, and correct it. A correct explanation becomes Pass; an incorrect explanation becomes Not Yet. \\
Not Yet & The response is incorrect, incomplete, or insufficient. & Retake the objective through the quiz-retake process. \\
\bottomrule
\end{tabularx}
\end{center}

* If you are unable to take a quiz when it is originally offered, you will need to use a retake for it (including filling out the Retake Qualification Form).

\textbf{Quiz correction \& retake procedure:}

You only need to address the objectives you did not pass—those marked Almost Pass or Not Yet. The steps are mostly the same for converting to a Pass and qualifying to retake; differences are marked \textbf{[convert]} or \textbf{[retake]}.

\begin{samepage}
\begin{enumerate}
\item[1.] \textbf{Prepare corrections} for each objective you missed:
\end{enumerate}
\end{samepage}

\begin{samepage}
\begin{itemize}
\item \textbf{[convert]} Be prepared to explain the correct answer, the reasoning behind it, and why you made a mistake. You may bring a copy of the printed quiz with your corrections, but it is not required.
\end{itemize}
\end{samepage}

\begin{samepage}
\begin{itemize}
\item \textbf{[retake]} Complete a \textbf{Quiz Retake Qualification Form}. Print the form and a blank copy of the quiz. Fill out the qualification form and the quiz corrections. For each question that was not fully correct, be prepared to explain the correct answer, the reasoning behind it, and why you made a mistake.
\end{itemize}
\end{samepage}

You may collaborate with others (e.g., during office hours or with a study group), but the corrections you submit must be your own, handwritten, and show all work.

\begin{samepage}
\begin{enumerate}
\item[2.] \textbf{Sign up for a correction meeting.} During the week after the quiz, sign up for a meeting with the course instructor or a teaching fellow.
\end{enumerate}
\end{samepage}

\begin{samepage}
\begin{enumerate}
\item[3.] \textbf{Attend your correction meeting.} At the meeting, you may have only your printed and completed corrections in front of you. The instructor will have your original submission open in Gradescope.
\end{enumerate}
\end{samepage}

\begin{samepage}
\begin{itemize}
\item \textbf{[retake]} Bring the completed retake qualification form as well.
\end{itemize}
\end{samepage}

\begin{samepage}
\begin{enumerate}
\item[4.] \textbf{Explain each incorrect objective.} You will be asked to explain what was wrong, why it was wrong, and what the correct answer is. Be prepared to answer follow-up questions.
\end{enumerate}
\end{samepage}

\begin{samepage}
\begin{enumerate}
\item[5.] \textbf{Outcome:}
\end{enumerate}
\end{samepage}

\begin{samepage}
\begin{itemize}
\item \textbf{[convert]} A correct explanation, including answers to follow-up questions, converts the objective to Pass. If you cannot explain it adequately, it becomes Not Yet. You will then have to retake the learning objective to obtain a Pass. In this case, you automatically qualify for the retake.
\end{itemize}
\end{samepage}

\begin{samepage}
\begin{itemize}
\item \textbf{[retake]} A satisfactory form and explanation qualify you to retake that objective. The instructor signs the form and records your qualification. This does not by itself make the objective Pass—you must still take the retake.
\end{itemize}
\end{samepage}

\begin{samepage}
\begin{enumerate}
\item[6.] \textbf{[retake] Complete the retake.} You will complete the retake questions at the end of the semester during the final exam period. Retakes are harder than the original quizzes, and each objective is graded Pass or Not Pass. A Not Pass result on a retake is final; there is no retake of a retake. There is no limit to the number of objectives for which you may qualify to retake.
\end{enumerate}
\end{samepage}

\section*{}
\textbf{LECTURE PARTICIPATION:}

Lecture participation is based on completion of in-class activities during the lecture.

\section*{}
\textbf{SKILLS LABS:}

The lab sections provide you the opportunity to work in groups to solve problems in probability theory and data analysis and learn to use the statistical programming language R using the RStudio software. About half of the labs will be “skills labs” and the other half will be dedicated to the lab projects (see below). During skills labs, you will have the chance to apply practical statistical concepts and deepen your understanding through hands-on work with R. Skills labs are evaluated based on submission of the in-lab worksheet and completion of the interactive tutorial. One completion of skills lab means submission of both the tutorial hash code and the worksheet for that lab.

\section*{}
\textbf{TUTORIALS:}

Interactive tutorials are provided to prepare students for the skills labs. Each tutorial includes guided explanations, R examples, and exercises. Each tutorial is available as an interactive lesson and must be completed before the associated lab.

\section*{}
\textbf{LAB PROJECTS:}

Students work in groups to complete two projects:

\begin{samepage}
\begin{itemize}
\item Project 1: video presentation
\item Project 2: written report
\end{itemize}
\end{samepage}

Each project is evaluated with the ESRN rubric:

\begin{center}
\small
\setlength{\tabcolsep}{4pt}
\begin{tabularx}{\textwidth}{lY}
\toprule
\textbf{Label} & \textbf{Meaning} \\
\midrule
E & Excellent or Exemplary \\
S & Satisfactory \\
R & Revisions Needed \\
N & Not Assessable \\
\bottomrule
\end{tabularx}
\end{center}

A grade of \textbf{S} receives full credit for the project while grades of \textbf{E} are considered when deciding between final letter grades. Projects receiving an \textbf{R} or \textbf{N} must be revised and resubmitted once. Any project that does not achieve an \textbf{S} or \textbf{E} by the end of the course will result in a reduction to your final course grade (see above).

\textbf{Peer Evaluation:}

For each project deliverable, every group member will submit a peer-assessment form. The peer-evaluation rubric includes the following categories: \textbf{Excellent}, \textbf{Satisfactory}, and \textbf{Unsatisfactory}. If you consistently receive \textbf{Unsatisfactory} ratings on the peer-evaluation form, the lab instructor will reach out to you. Repeated \textbf{Unsatisfactory} ratings in peer evaluation may result in a reduction of your final project grade.

\section*{}
\textbf{USE OF GENERATIVE AI:}

Generative artificial intelligence tools (GenAI) — software that creates new text, images, computer code, audio, video, and other content — have become widely available. Well-known examples include CoPilot, Gemini, ChatGPT, and Claude among others. \textbf{A complete policy on the use of Generative AI in this course is included on the course webpage}, including acceptable uses, unauthorized uses, and other considerations. In particular, no GenAI use is allowed for Quizzes. The goal of this policy is to allow you to use AI to enhance their learning, while limiting AI use that could provide an unfair advantage or detract from your learning. If you use generative AI tools to complete assignments in this course, in ways that are not explicitly authorized, the course instructors will apply the Boston University Code of Academic Integrity as appropriate to your specific case. If you are unsure about your use of AI, it is best to check in with the course instructor in advance. In addition, you must be wary of unintentional plagiarism or fabrication of data. Please act with integrity, for the sake of both your personal character and your academic record.

\section*{}
\textbf{ATTENDANCE:}

Students are expected to attend lectures, labs, and discussion sections unless they have a valid reason for being absent. When an absence is necessary, notify the lecture instructor, lab instructor, or discussion teaching fellow as soon as possible and ideally before the absence. Excused lecture absences are not penalized in the lecture participation requirement.

Boston University attendance policy: \url{https://www.bu.edu/academics/policies/attendance/}

\section*{}
\textbf{RECORDING:}

Lectures are recorded and shared with students who miss class for a valid reason. Every effort will be made to record lectures; however, technical difficulties may occasionally prevent a recording.

\section*{}
\textbf{ACADEMIC CONDUCT:}

All Boston University students are expected to maintain high standards of academic honesty and integrity. Students are responsible for understanding the Academic Conduct Code, including the ethical standards, rights, and responsibilities that apply to members of the university learning community. Cheating, plagiarism, and other forms of academic misconduct are addressed according to this policy. Penalties may range from failing an assignment or course to suspension or expulsion from the university.

Boston University Academic Conduct Code: \url{https://www.bu.edu/academics/policies/academic-conduct-code/}

\section*{}
\textbf{DISABILITY AND ACCESS SERVICES:}

Students with documented disabilities, including learning disabilities, may be entitled to accommodations that provide integrated and equal access to university programs. Accommodations may include additional testing time, staggered homework assignments, or note-taking assistance.

Students who believe they should receive accommodations should contact Disability and Access Services. The office can provide an accommodation letter to share with instructors; the letter does not state the reason for the accommodations. Students who already have a letter are encouraged to share it with the instructor as soon as possible.

Boston University Disability and Access Services: \url{https://www.bu.edu/disability/}

\section*{}
\textbf{EDUCATIONAL RESOURCE CENTER:}

The Educational Resource Center provides programs and services that help students strengthen their learning and develop a plan for academic success. Free peer tutoring may be available for this course, subject to availability.

Boston University Educational Resource Center: \url{https://www.bu.edu/erc/}

\section*{}
\textbf{WEEKLY PLAN (SUBJECT TO CHANGE):}

\begin{center}
\scriptsize
\setlength{\tabcolsep}{4pt}
\renewcommand{\arraystretch}{1.1}
\hyphenpenalty=10000 \exhyphenpenalty=10000
\begin{tabular}{|C{0.095\textwidth}|P{0.100\textwidth}|P{0.075\textwidth}|P{0.170\textwidth}|P{0.090\textwidth}|P{0.145\textwidth}|P{0.105\textwidth}|P{0.067\textwidth}|}
\hline
\textbf{Week} & \textbf{Reading (Chapter)} & \textbf{Module} & \textbf{Lecture Topics} & \textbf{Labs} & \begin{tabular}[t]{@{}c@{}}\textbf{Lab}\\\textbf{Deliverable}\end{tabular} & \textbf{Homework due Mon} & \textbf{Quiz} \\
\hline
1 (9/2) & 1; 2.1 & 1 &  & Lab 1 & Lab1 in-lab activity submission; Tutorial 1 &  &  \\
\hline
2 (9/7) & 2.2; 2.3 & 1 & Summarizing data & Lab 2 & Lab2 in-lab activity submission; Tutorial 2 & 1 &  \\
\hline
3 (9/14) & 3.1; 3.2–3.3; 3.4 & 1; 2 & Probability and RVs & Lab 3 & Lab3 in-lab activity submission; Tutorial 3 & 2 &  \\
\hline
4 (9/21) & 3.5; 4.1 & 2 & Continuous distributions & Project1-1 & Project1-1: Project 1 plan & 3 & Quiz 1 \\
\hline
5 (9/28) & 4.2; 4.3 & 2 & Normal distribution; Geometric distribution & Lab 4 & Project 1 Outline; Lab4 in-lab activity submission; Tutorial 4 & 4 &  \\
\hline
6 (10/5) & 4.5; 5.1 & 2; 3 & Binomial, Poisson distribution; Point estimates and sampling variability & Lab 5 & Lab5 in-lab activity submission; Tutorial 5 & 5 &  \\
\hline
7 (10/12)\newline Tues is\newline Mon schedule & 5.2; 5.3 & 3 & Point estimates and sampling variability; Confidence intervals & Project1-2 &  & 6 & Quiz 2 \\
\hline
8 (10/19) & 6.1; 6.2 & 3 & HTs; Inference for a single proportion & Project2-1 & Project2-1: Project1 Video & 7 &  \\
\hline
\end{tabular}
\end{center}
\newpage
\section*{}
\textbf{WEEKLY PLAN (CONTINUED):}
\begin{center}
\scriptsize
\setlength{\tabcolsep}{4pt}
\renewcommand{\arraystretch}{1.1}
\hyphenpenalty=10000 \exhyphenpenalty=10000
\begin{tabular}{|C{0.095\textwidth}|P{0.100\textwidth}|P{0.075\textwidth}|P{0.170\textwidth}|P{0.090\textwidth}|P{0.145\textwidth}|P{0.105\textwidth}|P{0.067\textwidth}|}
\hline
\textbf{Week} & \textbf{Reading (Chapter)} & \textbf{Module} & \textbf{Lecture Topics} & \textbf{Labs} & \begin{tabular}[t]{@{}c@{}}\textbf{Lab}\\\textbf{Deliverable}\end{tabular} & \textbf{Homework due Mon} & \textbf{Quiz} \\
\hline
9 (10/26) & 6.3; 6.4; 7.1 & 4 & Difference of two proportions, Chi-square GoF, two-way tables & Lab 6 & Lab6 in-lab activity submission; Tutorial 6; Project2 Outline & 8 &  \\
\hline
10 (11/2) & 7.2; 7.3 & 4 & Inference for one sample means, paired data, difference of two means & Project2-2 &  & 9 & Quiz 3 \\
\hline
11 (11/9) & 7.4; 7.5 & 4 & Power calculations, Issues with Hypothesis Testing, ANOVA & Project2-3 & Project2-3: Project2 Progress Report & 10 &  \\
\hline
12 (11/16) & 1 & 4; 5 & Frequentist and Bayesian Approaches to Statistical Modeling & Lab 7 & Project2 writeup; Lab7 in-lab activity submission; Tutorial 7 & 11 &  \\
\hline
13 (11/23)\newline No classes\newline Wed-Fri &  &  &  &  &  &  & Quiz 4 \\
\hline
14 (11/30) & 2; 3 & 5 & Bayesian modeling with discrete distributions &  & Project2 Resubmission & 12 &  \\
\hline
15 (12/7) &  & 5 & Special Topics &  &  &  & Quiz 5 \\
\hline
\end{tabular}
\end{center}

\end{document}
